% New introductory waypoint optimization exercise.
% Shared layout adapted from the instructor's ME 537 homework sources.
\documentclass[12pt,letterpaper]{article}
\usepackage{mathptmx}
\usepackage[T1]{fontenc}
\usepackage{xcolor,amsmath,amssymb,xurl,listings,enumitem,tikz}
\textwidth=6.5in
\textheight=9in
\oddsidemargin=0in
\evensidemargin=0in
\topmargin=0in
\headheight=0in
\headsep=0in
\parskip=0.08in
\parindent=0in
\setlist[enumerate]{itemsep=0.04in,topsep=0.04in}
\lstset{basicstyle=\ttfamily\footnotesize,breaklines=true,columns=fullflexible,
keepspaces=true,showstringspaces=false,aboveskip=0.04in,belowskip=0.04in}
\newcommand{\R}[2]{{}^{#1}\!R_{#2}}
\newcommand{\T}[2]{{}^{#1}\!T_{#2}}
\newcommand{\courseheader}[2]{%
\setlength{\fboxrule}{2pt}\setlength{\fboxsep}{2mm}%
\fcolorbox{black}{white}{\parbox{\dimexpr\linewidth-2\fboxsep-2\fboxrule\relax}{\centering\textbf{Introduction to Robotics Homework \##1 -- Draft}}}\par
\textbf{#2}\par}
\newcommand{\conventions}{Use right-handed rotations, column vectors, radians,
and meters. $\R{A}{B}$ maps coordinates from frame $B$ into frame $A$;
$\T{A}{B}$ also includes the origin of $B$ expressed in $A$.}
\newcommand{\submission}{\textbf{Submit:} labeled hand derivations and sketches,
your changed Python source files, the requested numerical/visual evidence,
and an AI-use statement (including ``none'' if applicable). Show your reasoning
before using code to check hand calculations. Public checks are intentionally
limited. The instructor supplies the deadline and submission destination.}

\begin{document}
\courseheader{7}{Waypoint optimization and path validation}
\textbf{Schedule:} second part of the planning unit. Build on HW06.
These are geometric paths: the waypoint index is not time.

\begin{enumerate}
\item \textbf{Formulate and calculate a path objective.}
For waypoints $q_0,\ldots,q_N$, consider the cost
\[
C(q)=\sum_{i=1}^{N-1}\|q_{i+1}-2q_i+q_{i-1}\|^2.
\]
\begin{enumerate}
\item Identify the decision variables if $q_0$ and $q_N$ are fixed. State
the endpoint, joint-limit, and collision requirements for an arm path.
Which requirements are constraints, and what is the initial guess?
\item For $q_0=[0,0]$, $q_1=[1,1]$, $q_2=[2,0]$, and $q_3=[3,0]$,
calculate $C(q)$ and the following path length by hand:
\[
L=\sum_{i=0}^{2}\|q_{i+1}-q_i\|.
\]
\item Explain why $C$ measures changes between waypoint steps. Why does it
not establish physical joint acceleration without a timing model?
\end{enumerate}

\item \textbf{Implement two short path metrics.}\par
The assigned file is:
\begin{lstlisting}
src/intro_robotics/waypoint_metrics.py
\end{lstlisting}
Complete \texttt{path\_length} and \texttt{smoothness\_cost}. The input checks are
provided in an imported utility. Keep the supplied function signatures.
\begin{enumerate}
\item Use consecutive differences for length and second differences for the
cost. Check your functions on problem 1(b) and add one example of your own.
\item From \path{/workspace} in the course container, run the limited checks:
\begin{lstlisting}
python -m homework.hw07.hw07_optimization
\end{lstlisting}
For a two-revolute-joint path, what units does your length use? Why is it not
the Cartesian distance traveled by the tool?
\end{enumerate}
\end{enumerate}

\newpage
\begin{enumerate}[start=3]
\item \textbf{Compare two initial paths for a circular base.}
The provided scene has a circular obstacle of radius $0.30$ m and a base of
radius $0.06$ m. The supplied optimizer uses a smoothness cost plus a soft
penalty for waypoint proximity to the expanded obstacle.
\begin{enumerate}
\item Use \texttt{base\_problem()} to obtain an RRT seed and a straight seed.
Run the provided optimizer on each with the same settings. Report initial/final
objective, path length, and whether each edge is valid at resolution $0.005$ m.
Use \texttt{sampled\_path\_is\_free} for this check.
See \path{homework/hw07/README.md} for a short starter script.
\item Explain the difference between a soft penalty and a hard collision
constraint. Does a lower objective alone prove that a path is collision-free?
Explain the observed result from each seed.
\item Change the obstacle weight for one run and compare smoothness and
validity. Total objective values with different weights are different
objectives; do not compare them as if they were the same quantity.
\end{enumerate}

\item \textbf{Refine an arm path.}
Use \texttt{arm\_problem()} for the HW06 arm. The nominal link radius is
$0.015$ m; the provided planning scene uses $0.030$ m to add clearance.
The provided \texttt{refine\_feasible\_path} keeps endpoints fixed and
accepts updates only when all sampled edges pass the selected checker.
\begin{enumerate}
\item Refine the provided RRT seed and report its initial/final smoothness,
joint-space length, optimizer time, and returned validity flag. Independently
check every returned edge against the nominal scene at resolution $0.005$ rad.
\item Try the direct joint-space seed. What does the optimizer report about
its initial validity? Explain why it cannot conclude that no valid route exists.
\item Explain the role of the extra clearance and why even the finer sampled
check is not a proof of continuous collision freedom.
\end{enumerate}

\item \textbf{Visualize and interpret.}
\begin{lstlisting}
python -m homework.hw07.hw07_optimization --show
python -m homework.hw07.hw07_optimization --save tmp/hw07.png
\end{lstlisting}
Save a view and explain one reason to combine RRT-Connect with local
optimization. State one limitation of the models or comparison. The supplied
optimizers are simple teaching examples; they are not CHOMP implementations.
\end{enumerate}

\submission
\end{document}
